% Source: https://artofproblemsolving.com/community/c3h2314921_modular_arithmetic
% Stable identity: aops-topic-2314921-q01
\textbf{(a)} Since $\gcd(5,2)=1$, $5$ has an inverse modulo $2$: $5\cdot1=5\equiv1\pmod2$, so $5^{-1}\equiv1$. Then $\dfrac15\equiv1\cdot1=1$ and $\dfrac95\equiv9\cdot1\equiv1\pmod2$, and
\[\dfrac15+\dfrac95\equiv1+1=2\equiv0\pmod2.\]
So the statement makes sense and is $\boxed{\text{true}}$.

\textbf{(b)} As rational numbers, $\dfrac15+\dfrac95=\dfrac{10}{5}=2\equiv0\pmod2$, the same result. This is no accident: division by a number coprime to the modulus is multiplication by its inverse, which is compatible with the usual rules of fractions.

\textbf{(c)} $\dfrac12$ modulo $2$ would be a number $y$ with $2y\equiv1\pmod2$. But $2y$ is always even, so $2y\equiv0\not\equiv1\pmod2$. Since $\gcd(2,2)=2\ne1$, the number $2$ has no inverse modulo $2$, and $\dfrac12$ is undefined.
